\documentclass[oneside,a4paper,14pt]{extarticle}
\usepackage[a4paper,letterpaper,top=20mm,bottom=20mm,left=20mm,right=10mm]{geometry}
\usepackage[russian]{babel}
\usepackage{indentfirst}
\usepackage{graphicx}
\usepackage{caption}
\usepackage{titlesec}
\usepackage{minted, fancyvrb}
\usepackage{hyperref}
\usepackage{enumitem}
\usepackage{float}
\DeclareUnicodeCharacter{2500}{-}
\DeclareUnicodeCharacter{1F3CE}{} % 🏎
\DeclareUnicodeCharacter{1F3C6}{} % 🏆
\DeclareUnicodeCharacter{1F947}{} % 🥇
\DeclareUnicodeCharacter{26A1}{}  % ⚡
\DeclareUnicodeCharacter{1F4CB}{} % 📋
\DeclareUnicodeCharacter{1F3D7}{} % 🏗
\DeclareUnicodeCharacter{1F3C1}{} % 🏁
\DeclareUnicodeCharacter{1F468}{} % 🧑
\DeclareUnicodeCharacter{200D}{}  % ZWJ
\DeclareUnicodeCharacter{1F680}{} % 🚀
\DeclareUnicodeCharacter{270E}{}  % ✎

% Code styling
\setminted{style = rainbow_dash, fontsize = \small, breaklines}

% Header styles
\titleformat{\section}{\normalsize\bfseries}{\thesection}{1em}{}
\titleformat{\subsection}{\normalsize\bfseries}{\thesubsection}{1em}{}
\titleformat{\subsubsection}{\normalsize\bfseries}{\thesubsubsection}{1em}{}

% Spacing
\renewcommand\baselinestretch{1.33}
\setlength{\parindent}{1.25cm}

% List styles
\setlist[enumerate]{
  left=\parindent,
  label=\arabic*.,
  itemsep=0pt,
  topsep=5pt,
  partopsep=0pt,
  parsep=0pt
}
\setlist[itemize]{
  left=\parindent,
  itemsep=0pt,
  topsep=5pt,
  partopsep=0pt,
  parsep=0pt
}

% Hyperlinks
\hypersetup{
  colorlinks=true,
  linkcolor=black,
  urlcolor=blue,
  pdfborder={0 0 0},
  pdftitle={Система управления гоночными командами — Лабораторная работа 4},
  pdfauthor={Черкасов А.А.}
}

\begin{document}

\newpage
\thispagestyle{empty}
\begin{center}
  МИНИСТЕРСТВО НАУКИ И ВЫСШЕГО ОБРАЗОВАНИЯ РОССИЙСКОЙ ФЕДЕРАЦИИ ФЕДЕРАЛЬНОЕ ГОСУДАРСТВЕННОЕ БЮДЖЕТНОЕ ОБРАЗОВАТЕЛЬНОЕ УЧРЕЖДЕНИЕ ВЫСШЕГО ОБРАЗОВАНИЯ\\
  «ВЯТСКИЙ ГОСУДАРСТВЕННЫЙ УНИВЕРСИТЕТ»\\
  Институт математики и информационных систем\\
  Факультет автоматики и вычислительной техники\\
  Кафедра электронных вычислительных машин
\end{center}
\vspace{10mm}

\hfill
\begin{tabular}{l}
  \footnotesize Дата сдачи на проверку:                                          \\
  \footnotesize <<\rule[-1mm]{5mm}{0.10mm}\/>>\ \rule[-1mm]{20mm}{0.10mm}\ 2026 г. \\
  \footnotesize Проверено:                                                       \\
  \footnotesize <<\rule[-1mm]{5mm}{0.10mm}\/>>\ \rule[-1mm]{20mm}{0.10mm}\ 2026 г. \\
\end{tabular}
\vfill

\begin{center}
  Разработка распределённой веб-системы Formula Racing Management System на Go и Python с базой данных PostgreSQL. \\
  Отчёт по лабораторной работе №4\\
  по дисциплине\\
  <<Объектно-ориентированное программирование>>\\
\end{center}
\vspace{25mm}
\noindent
\begin{tabular}{ll}
  Разработал студент гр. ИВТб-2301-05-00 & \hspace{18mm}\rule[-1mm]{30mm}{0.10mm}\,/Черкасов А. А./ \\
                                         & \hspace{25.5mm}\footnotesize(подпись)                    \\
  Преподаватель                  & \hspace{18mm}\rule[-1mm]{30mm}{0.10mm}\,/Шмакова Н. А./ \\
                                         & \hspace{25.5mm}\footnotesize(подпись)                    \\
\end{tabular}

\noindent
\begin{tabular}{lp{58mm}r}
  Работа защищена &  & \hspace{13mm}<<\rule[-1mm]{5mm}{0.10mm}\/>>\ \rule[-1mm]{30mm}{0.10mm}\ 2026 г.
\end{tabular}
\vfill

\begin{center}
  Киров\\
  2026
\end{center}

% ──────────────────────────────────────────────────────────────────────────────
\newpage\thispagestyle{plain}

\section*{Цели лабораторной работы}
\begin{itemize}
  \item[$-$] разработать распределённое веб-приложение на базе микросервисной архитектуры;
  \item[$-$] реализовать бизнес-логику и веб-интерфейс на языке Go;
  \item[$-$] реализовать слой доступа к данным на Python через ORM SQLAlchemy, подключившись к существующей схеме базы данных из предыдущих лабораторных работ (PostgreSQL);
  \item[$-$] обеспечить взаимодействие между сервисами через HTTP (REST API).
\end{itemize}

\section*{Задание}
\begin{enumerate}
  \item Использовать базу данных из Лабораторной работы №3 (PostgreSQL) без модификации структуры таблиц.
  \item Написать слой доступа к данным на Python с использованием SQLAlchemy и паттерна <<Репозиторий>> (CRUD).
  \item Написать бизнес-логику на Go, реализующую: интерфейс пользователя (HTML/CSS), фильтрацию данных, агрегацию (сумма/среднее), валидацию форм, кэширование In-Memory с TTL, проксирование и методы из ЛР №1.
  \item Реализовать CRUD-интерфейс для всех объектов: Просмотр, Редактирование, Удаление, Добавление.
  \item Обеспечить выпадающие списки (Foreign Keys) и отображение связанных сущностей при просмотре объектов (например, Команды и Серии).
\end{enumerate}

% ──────────────────────────────────────────────────────────────────────────────
\clearpage
\section*{Реализация приложения}

\subsection*{Архитектура распределённой системы}
Проект <<Formula Racing>> реализован в виде трёхзвенной распределённой архитектуры:
\begin{enumerate}
    \item \textbf{Клиент (Браузер):} Получает отрендеренные \texttt{HTML/CSS/JS} страницы с сервера Go. Тематическое цветовое оформление реализовано с использованием набора цвета Catppuccin Mocha.
    \item \textbf{BFF-Сервер (Backend-for-Frontend, Go):} Точка входа для браузера. Выполняет функции рендеринга шаблонов, кэширования запросов к Python-серверу на протяжении 5 минут (in-memory TTLCache), маршрутизации REST-запросов и вычисления бизнес-логики (обработка списков пилотов).
    \item \textbf{Data-API Сервер (Python/Flask):} Инкапсулирует подключение к существующей базе данных PostgreSQL. Обрабатывает REST-запросы от Go и выполняет SQL-команды используя паттерн ООП Repository.
\end{enumerate}

\subsection*{Моделирование репозитория на Python}
Доступ к данным обеспечен через ORM-библиотеку SQLAlchemy. Все сущности строго соответствуют схеме базы данных PostgreSQL (созданной Flyway в ЛР №3): \texttt{racing\_series}, \texttt{racing\_teams}, \texttt{racing\_drivers}.

Для реализации внешних ключей и связей применены \texttt{ForeignKey(..., ondelete='CASCADE')} и \texttt{relationship}:
\begin{itemize}
  \item \textbf{В классе RacingDriver:} поле \texttt{team\_id = Column(BigInteger, ForeignKey("racing\_teams.id"))} связывает пилота с его командой, \texttt{relationship} позволяет подгружать связанные сущности через механизмы eager-loading.
  \item \textbf{В классе RacingTeam:} поле \texttt{series\_id} указывает на гоночный чемпионат (\texttt{RacingSeries}).
\end{itemize}

Слой API представляет собой приложение на основе Flask, объявляющее роуты (роутинг \texttt{@app.route}) к функциям-репозиториям, скрывающим сложную логику работы с сессиями.

\subsection*{Бизнес-логика и интерфейс на Go}
Клиентская логика и веб-интерфейс разработаны в пакете Go при помощи стандартных библиотек \texttt{net/http} и \texttt{html/template}.

\textbf{Функции на стороне Go включают:}
\begin{itemize}
    \item[$-$] \textbf{ООП-методы (ЛР №1):} Функции \texttt{CalculateF1Points}, \texttt{GetWinRate} написаны на Go и вызываются прямо внутри HTML-шаблонов благодаря \texttt{template.FuncMap}.
    \item[$-$] \textbf{Агрегация и фильтрация:} Функция \texttt{ComputeStats} обходит массивы структур и высчитывает агрегированные данные: среднее количество очков, максимальное количество очков у лидирующего пилота. Функция \texttt{SortDrivers} позволяет клиенту динамически отсортировать результат.
    \item[$-$] \textbf{Валидация перед отправкой:} Функции проверки бизнес-правил (\texttt{ValidateTeam}, \texttt{ValidateDriver}) исполняются на сервере Go перед отправкой запроса на сервер Python. Таким образом Data-сервер изолирован от "грязных" данных.
    \item[$-$] \textbf{In-Memory кэширование:} Реализован модуль \texttt{cache.go} с использованием RW-блокировок (\texttt{sync.RWMutex}) и встроенного сборщика мусора на независимой горутине. При частых GET-запросах, результат возвращается мгновенно, минуя Python-сервер (TTL 5 минут).
\end{itemize}

Для добавления или изменения Команд/Пилотов используются HTML-формы с выпадающими списками (загружающими связанные чемпионаты или команды через подгруженные Foreign Keys с Python API), реализуя полноценный CRUD-цикл.

\section*{Вывод}
В результате выполнения лабораторной работы была успешно внедрена распределённая клиент-серверная микросервисная архитектура без изменения исходной схемы базы данных. Перенос операций хранения на Python с использованием ORM-репозиториев в связке с легковесным, кеширующим сервером отображения на Go позволил получить масштабируемую, стабильную и независимую систему с удобным графическим Catppuccin веб-интерфейсом.

\newpage

\section*{Приложение А1. Go: Выделенные предметные модели (models.go)}
\inputminted{go}{code/go/models.go}

\section*{Приложение А2. Go: Бизнес-логика и агрегация (business.go)}
\inputminted{go}{code/go/business.go}

\section*{Приложение А3. Go: Модуль In-Memory кэширования (cache.go)}
\inputminted{go}{code/go/cache.go}

\section*{Приложение А4. Go: API-клиент (api\_client.go)}
\inputminted{go}{code/go/api_client.go}

\section*{Приложение А5. Go: Основные HTTP обработчики (handlers.go)}
\inputminted{go}{code/go/handlers.go}

\section*{Приложение А6. Go: Обертка инициализации (main.go)}
\inputminted{go}{code/go/main.go}

\section*{Приложение В1. Python: Модели БД SQLAlchemy (models.py)}
\inputminted{python}{code/python/models.py}

\section*{Приложение В2. Python: CRUD-репозиторий (repository.py)}
\inputminted{python}{code/python/repository.py}

\section*{Приложение В3. Python: Flask REST API (app.py)}
\inputminted{python}{code/python/app.py}

\section*{Приложение С1. Веб-интерфейс: Главная страница (templates/index.html)}
\inputminted{html}{code/go/templates/index.html}

\section*{Приложение С2. Веб-интерфейс: Список команд (templates/teams.html)}
\inputminted{html}{code/go/templates/teams.html}

\section*{Приложение С3. Веб-интерфейс: Список пилотов (templates/drivers.html)}
\inputminted{html}{code/go/templates/drivers.html}

\section*{Приложение С4. Веб-интерфейс: Стили (static/style.css)}
\inputminted{css}{code/go/static/style.css}

\end{document}
